\scanpage{339}
\mspage{78}
% archive sheet number: 338
formés d'un hom $\unsure{w|u}:c_1\to c_2$ et $w':c'_1\mathrel{\unsure{-|\to}}c'_2$ dans $\underline{\mathrm{Ch}}_{\Sigma,\tau}(a,b)$ et $\underline{\mathrm{Ch}}_{\Sigma,\tau'}(b,c)$ respectivement, et si $w'',w,w'$ se correspondent ainsi, on posera
\[
w''=w*w'.
\]

(Il y a une transitivité évidente à vérifier pour voir que ça définit un \uline{foncteur} $\unsure{(*)}$. On peut exprimer la situation à l'aide du diagramme
\diagram{\begin{tikzcd}[ampersand replacement=\&,row sep=large,column sep=small]
\& \underline{\mathrm{Ty}}\times \underline{\mathrm{Ty}}\arrow[rr,"{\Pi\;(\text{composition})}"]\arrow[dl,"p_1"']\arrow[dr,"p_2"] \& \& \underline{\mathrm{Ty}}\\
\underline{\mathrm{Ty}} \& \& \underline{\mathrm{Ty}} \& {,}
\end{tikzcd}}
où $p_1,p_2$ sont les deux projections. On considère sur $\underline{\mathrm{Ty}}\times \underline{\mathrm{Ty}}$ les trois catégories fibrées
\[
p_1^*(\underline{\mathrm{CH}}_{\Sigma,\tau}(a,b)),\quad
p_2^*(\underline{\mathrm{CH}}_{\Sigma,\tau'}(b,\unsure{c})),\quad
\Pi^*(\underline{\mathrm{Ch}}_{\Sigma,\tau\tau'}(a,c))
\]
dont les fibres en $(\tau,\tau')$ sont respectivement
\[
\underline{\mathrm{Ch}}_{\Sigma,\tau}(a,b),\quad
\underline{\mathrm{Ch}}_{\Sigma,\tau'}(b,c),\quad
\underline{\mathrm{Ch}}_{\Sigma,\tau\tau'}(\unsure{\text{à}|a},c).
\]

(On \unsure{a} écrit $\underline{\mathrm{CH}}$ au lieu de $\underline{\mathrm{Ch}}$, pour indiquer que l'on considère les catégories envisagées comme des catégories (fibrées) \uline{au dessus de} $(\underline{\mathrm{Ty}})$.)
% slanted
\margin{Mais il s'agit pour $p_1^*,p_2^*$ et $\Pi^*$ de produits cartésiens ordinaires vus dans $\mathrm{Cat}$, il n'y a pas \unsure{à avoir|à voir|raison} pour \unsure{cela} à écrire des ``$\mathrm{CH}$''\dots}

Au dessus de $\underline{\mathrm{Ty}}\times \underline{\mathrm{Ty}}$, on a un \deleted{\unsure{foncteur}} (cartésien) de catégories fibrées \deleted{sur $\underline{\mathrm{Ty}}\times$}
\[
p_{\unsure{1|2}}^*\underline{\mathrm{CH}}_{\tau}(a,b))\underset{\underline{\mathrm{Ty}}\times \underline{\mathrm{Ty}}}{\times}
p_2^*(\underline{\mathrm{CH}}_{\tau}(b,c))\longrightarrow\Pi^*(\underline{\mathrm{CH}}_{\tau}(a,c)),\tag*{$(**)$}
\]
